\documentclass[10pt,a4paper]{amsart}
\usepackage[margin=19mm]{geometry}
\usepackage{amsmath,amssymb,mathtools}
\usepackage[scr=rsfs,cal=euler]{mathalfa}
\usepackage{xfrac}
\usepackage[unicode,hidelinks,bookmarks=false]{hyperref}
\newcommand{\ie}{i.e.\ }
\newcommand{\cf}{cf.\ }
\newcommand{\resp}{respectively\ }
\newcommand{\Subsection}{\subsection}
\providecommand{\nobreakdash}{\nobreak\hbox{-}}
\setlength{\emergencystretch}{2em}
\multlinegap0pt
% ---------------------------------------------------------------------------
% Editorial AMS wrapper prelude. The theorem environments and the mathematical macros below are
% transcribed from the paper's own preamble (cached-originals/source0.tex lines 41--58); the AMS amsart
% class, the named format packages of the pinned TeX Live runtime and this editorial formatting are not
% part of the author's text. No mathematics is changed.
% ---------------------------------------------------------------------------
\usepackage{tikz-cd}
\newtheorem{theorem}{Theorem}[section]
\newtheorem{proposition}[theorem]{Proposition}
\newtheorem{definition}[theorem]{Definition}
\newtheorem{corollary}{Corollary}[theorem]
\newtheorem{note}{Note}[section]
\newtheorem{lemma}[theorem]{Lemma}
\newtheorem*{note*}{Note}
\newtheorem*{claim*}{Claim}
\newtheorem*{lemma*}{Lemma}
\newtheorem*{examples*}{example}
\newtheorem*{example*}{example}
\newtheorem*{corollary*}{Corollary}
\newtheorem*{theorem*}{Theorem}
\newtheorem*{Reminder*}{Reminder}
\newtheorem*{justification*}{Justification}
\newtheorem*{Notation Convention*}{Notation Convention}
\newcommand{\la}[1]{\mathrm{#1}}
\newcommand{\ca}[1]{\mathcal{#1}}
\newcommand{\bb}[1]{\mathbb{#1}}
\newcommand{\cat}[1]{\mathsf{#1}}
\newcommand{\go}[1]{\mathfrak{#1}}
\newcommand{\lu}[1]{\mathscr{#1}}
\newcommand{\Hom}{\mathrm{Hom}}
\newcommand{\forallw}{\forall_{\omega}}
\newcommand{\dash}{\text{-}}


\begin{document}
\title{Generalised ultracategories and conceptual completeness of geometric logic}
\author{Ali Hamad}
\address{University of Ottawa}
\email{ahama099@uottawa.ca}
\date{Source: arXiv:2507.07922v4; distributed under CC BY 4.0}
\maketitle
\noindent\emph{Source, version and licence.} \emph{Generalised ultracategories and conceptual completeness of geometric logic}, by Ali Hamad. The source of this editable unit is the e-print of \textbf{version v4} of arXiv:2507.07922, https://arxiv.org/abs/2507.07922v4, whose material is distributed under the Creative Commons Attribution 4.0 International licence, http://creativecommons.org/licenses/by/4.0/, as recorded in the arXiv licence metadata for this paper and shown on that abstract page. The whole of the body below is version v4, the newest version of the paper. Full rights notice: licenses/arxiv-2507.07922-CC-BY-4.0.txt.\par\smallskip
\noindent\textit{Editorial scope.} Counted unit: the article's Lemma in Section 4, the statement carried by the source environment \texttt{lemma} with source0.tex lines 1878--1902, which says that for two logical topological groupoids $G(P,I)$ and $G(P^{\prime},I^{\prime})$ and a morphism $f$ in the image of the category $\ca{P}$ under $G$, the natural comparison square between them commutes up to natural isomorphism; it is printed with the number the article gives it. It is delivered together with the author's own complete proof as the single primary proof body, source0.tex lines 1904--1919 (16 lines): the \texttt{proof} environment that immediately follows the statement, that takes a sheaf $B$ in the topos $E$, reads it as a sheaf with action on both groupoids along the two geometric morphisms, compares the étale spaces $p^*(B)$ and ${p^{\prime}}^*(B)$ over $G_0$ and $G^{\prime}_0$, and checks that the pullback topology agrees with the topology of the étale space on the level of basic open sets. No proof marker is added and no mathematics is written, repaired or re-derived; the author's notation and the printed slips of the source are kept literal. Necessary context is retained uncounted, in source order: source0.tex lines 1209--1212, the article's Main theorem on the conceptual completeness of geometric logic for toposes with enough points, and lines 1820--1822, the third equivalence theorem of its proof, inside whose proof the counted lemma sits. The article's own numbering is reproduced by explicit counter settings: the counter \texttt{theorem}, declared by \texttt{\textbackslash newtheorem\{theorem\}\{Theorem\}[section]} and shared by proposition, definition, corollary and lemma, is set before each retained passage, so the two retained theorem statements print with the numbers the article gives them. No \texttt{\textbackslash cite} occurs in any retained passage: the closure of the retained passages over references and citations was computed, it contains no reference and no citation at all, so this unit delivers no bibliography entry and the article's own bibliography data file \texttt{bib.bib} is neither spliced into the bundled source nor redistributed. The retained passages contain no figure environment and no graphics-inclusion command, so no artwork is redistributed; the commutative diagrams of the counted statement and proof are drawn with the named \texttt{tikz-cd} package of the pinned TeX Live runtime, exactly as the source does. The source's own class is the standard \texttt{article} class; the e-print ships the authors' own style file \texttt{quiver.sty}, which is not shipped, loaded or redistributed here (the mathematical macros the retained passages need are transcribed from the paper's own preamble instead), and no publisher class is shipped, used or redistributed. Every declared body of this unit is an exact, unmodified span of source0.tex: no body is a bounded passage comparison, \texttt{omit} was not used anywhere, and consequently no fidelity receipt is asserted in delivery-evidence.json. The complete concatenated original source is bundled as sources/181378/source0.tex. Omitted from this editable unit and therefore absent from it are source0.tex lines 1--1208; 1213--1819; 1823--1877; 1903--1903; 1920--3675, that is the article's own preamble and title block, the introduction, the preliminaries on generalised ultracategories, the construction of the category $\ca{P}$ and of the functor $G$ with the two cited references it carries, the first and second equivalence theorems of the Main theorem with their proofs, and the article's bibliography data, all of which remain present in the bundled source but are not delivered here. The AMS \texttt{amsart} wrapper, the named format packages of the pinned TeX Live runtime and this editorial source, licence and scope paragraph are editorial formatting only.\par\smallskip
\setcounter{section}{7}
\setcounter{equation}{0}
\setcounter{corollary}{0}
\setcounter{note}{0}
\setcounter{theorem}{0}
\begin{theorem} 
\label{Main theorem}
    Let $E$ and $E^{'}$ be two toposes with enough points, then there is an equivalence of categories between $\mathrm{Lult}(M_{E}, M_{E^{'}})$ and $\mathrm{Geom}(E,E^{'})$.
\end{theorem}

\setcounter{section}{7}
\setcounter{equation}{0}
\setcounter{corollary}{0}
\setcounter{note}{0}
\setcounter{theorem}{5}
\begin{theorem}
    Let $E$ and $E^{'}$ be two toposes with enough points, then there is an equivalence of categories between $\la{Geom}(E,E^{'})$ and $\la{ClovenCart}(\cat{Top}//E, \cat{Top}//E^{'})$.\label{third equivalence theorem}
\end{theorem}

\setcounter{section}{7}
\setcounter{equation}{0}
\setcounter{corollary}{0}
\setcounter{note}{0}
\setcounter{theorem}{8}
\begin{lemma} Suppose that we have two topological groupoids $G(P,I)$ and $G(P^{'},I')$, and a morphism $f$ between them which is in the image of the category $\ca{P}$ by $G$, then the following diagram commutes up to natural isomorphism:


\[
% https://q.uiver.app/#q=WzAsNSxbMCwyLCJHXzAiXSxbMCwwLCJHXzEiXSxbMiw0LCJFIl0sWzQsMiwiR157J31fMCJdLFs0LDAsIkdeeyd9XzEiXSxbMSwwLCIiLDAseyJsYWJlbF9wb3NpdGlvbiI6NzAsIm9mZnNldCI6M31dLFsxLDAsIiIsMix7Im9mZnNldCI6LTN9XSxbMCwyLCJwIiwxXSxbMCwzLCJmXzAiLDFdLFs0LDMsIiIsMix7Im9mZnNldCI6LTN9XSxbNCwzLCIiLDAseyJvZmZzZXQiOjN9XSxbMSw0LCJmXzEiLDFdLFszLDIsInBeeyd9IiwxXSxbNywxMiwiXFxzaW1lcSIsMSx7InNob3J0ZW4iOnsic291cmNlIjoyMCwidGFyZ2V0IjoyMH19XV0=
\begin{tikzcd}
	{G_1} &&&& {G^{'}_1} \\
	\\
	{G_0} &&&& {G^{'}_0} \\
	\\
	&& E
	\arrow["{f_1}"{description}, from=1-1, to=1-5]
	\arrow[shift right=3, from=1-1, to=3-1]
	\arrow[shift left=3, from=1-1, to=3-1]
	\arrow[shift left=3, from=1-5, to=3-5]
	\arrow[shift right=3, from=1-5, to=3-5]
	\arrow["{f_0}"{description}, from=3-1, to=3-5]
	\arrow[""{name=0, anchor=center, inner sep=0}, "p"{description}, from=3-1, to=5-3]
	\arrow[""{name=1, anchor=center, inner sep=0}, "{p^{'}}"{description}, from=3-5, to=5-3]
	\arrow["\simeq"{description}, shorten <=13pt, shorten >=13pt, Rightarrow, from=0, to=1]
\end{tikzcd}
\]


\end{lemma}

\setcounter{section}{7}
\setcounter{equation}{0}
\setcounter{corollary}{0}
\setcounter{note}{0}
\setcounter{theorem}{9}
\begin{proof}
    

To see this, let $B$ be a sheaf in $E$, and let $p$ and $p^{'}$ be the respective geometric morphisms from  $G(P,I)$ and $G^{'}(I^{'},P^{'})$, then $B$ can be regarded as a sheaf on $G_0$ (image by $p^*$) with an action of $G_1$, or as a sheaf on $G^{'}_0$ (image by ${p^{'}}^*$) with an action of $G^{'}_1$, now let us regard the following diagram: 

\[\begin{tikzcd}
	{{p^{*}(B)} } && {{{p'}^{*}(B)}} \\
	\\
	{G_0} && {G'_0}
	\arrow["{\pi_1}"', from=1-1, to=3-1]
	\arrow["{\pi_2}", from=1-3, to=3-3]
	\arrow["{f_0}"', from=3-1, to=3-3]
\end{tikzcd}\]
Our goal is to show that the topology of the étale space $p^*(B)$ is the pullback of that of ${p'}^*(B)$, first it's easy to check that as sets, the set  of $p^*(B)$  is the set of the pullback (up to iso). It remains to check that it has the pullback topology, but one can easily notice that the basic open sets in both coincide.

\end{proof}

\end{document}
